% !TEX root = ../main.tex

\chapter{（这是笔记名）测试笔记}
\section{（这是第一章）夜色名为温柔}

\subsection{（这是第一节）天空与你之间}

\subsubsection{（这是第一小节）萤和球棒侠}

\textbf{接下来用于测试基本文本编辑，包括各种有序列表、无序列表、行内公式以及字体颜色等：}

\begin{itemize}
	\item 开路电压是指没有外部负载连接时，热电发电装置两端的电压。根据Seebeck效应，开路电压与温差成正比，$V_{\text{oc}} = \alpha \Delta T$。
	
	\begin{itemize}
		\item \textcolor{fred}{通过实验测量，展示即使是优化过的热电装置，其效率也受到卡诺效率的限制，这直接体现了热力学第二定律。}
		\begin{itemize}
			\item \textcolor{fred}{通过实验测量，展示即使是优化过的热电装置，其效率也受到卡诺效率的限制，这直接体现了热力学第二定律。}
			\item 讨论如何通过实验设计来逼近卡诺效率，包括使用最佳材料、最佳温差和最佳负载条件。
		\end{itemize}			
		\item 讨论如何通过实验设计来逼近卡诺效率，包括使用最佳材料、最佳温差和最佳负载条件。
	\end{itemize}			
	
	\item 当热电发电装置连接到负载时，输出电压会由于内阻的存在而下降。负载电压$V_{\text{L}} = \frac{\alpha \Delta T \cdot R_{\text{L}}}{R_{\text{L}} + R_{\text{in}}}$，而输出功率是负载上消耗的功率$P_{\text{out}} = \frac{(\alpha \Delta T)^2 \cdot R_{\text{L}}}{(R_{\text{L}} + R_{\text{in}})^2}$。
	
	\item 当负载电阻等于内阻时，热电发电装置输出的功率达到最大$P_{\text{max}} = \frac{(\alpha \Delta T)^2}{4 R_{\text{in}}}$。				
\end{itemize}

\begin{enumerate}
	\item 各元件性能测量
	\begin{enumerate}
		\item \textcolor{fred}{Seebeck效应测量}：通过改变温差并测量开路电压来研究Seebeck效应，从而确定$\alpha$；
		\begin{enumerate}
			\item \textcolor{fred}{热电堆集成——电加热器与Seebeck元件，测控部分——PID控温（热端、冷端），输出电路部分（输出功率测量）。}
			\begin{enumerate}
				\item \textcolor{fred}{加热功率测控}——电流表电压表实时测控（含于PID系统中，编写相关程序）；
				\item \textcolor{fred}{输出功率测控}——电流表电压表实时测控（考虑是否编程）；
				\item 测量最大输出功率——改变输出电路的负载；
				\item 测量最大输出效率——改变温差与负载，优化其它部分；
				\item 探究如何测量热端的损失。
			\end{enumerate}
		\end{enumerate}
		\item 器件内阻的测量与影响：内阻对热电转换效率有重要影响，了解并优化内阻对提高热机性能是必要的；
		\item 考虑测量电加热器的加热功率（以及其它元件性能，如Peltier效应）。
	\end{enumerate}
	
	\item \textcolor{fred}{搭建热机}
	
	\item 负载与效率测量
	
	
	\item 热力学第二定律的展示
	\begin{itemize}
		\item \textcolor{fred}{通过实验测量，展示即使是优化过的热电装置，其效率也受到卡诺效率的限制，这直接体现了热力学第二定律。}
		\item 讨论如何通过实验设计来逼近卡诺效率，包括使用最佳材料、最佳温差和最佳负载条件。
	\end{itemize}			
	
	\item 结论与进一步的探索
	\begin{itemize}
		\item 比较热机效率与理论卡诺效率的差异，讨论可能的优化途径和技术挑战；
		\item 探究各部分如何\textcolor{fred}{优化}（热电堆——减小散热，增大温差，冷热端优化；测控——优化控温算法，改进实时测量程序；输出电路部分——减小散热，考虑Thomson效应的影响；理论建模——估算散热进一步修正）。
	\end{itemize}
\end{enumerate}

\textbf{接下来用于测试图片、表格以及带编号公式，并且考虑它们的引用：}

\textbf{这是带编号公式：}

\begin{equation}
	V_{\text{oc}} = \alpha \Delta T
	\label{eq:test}
\end{equation}

来看看\cref{eq:test}。

\textbf{这是表格：}

\begin{table}[h]
	\centering
	\caption{高级控温算法的比较}
	\label{tab:tab1}
	\begin{tabular}{|c|c|c|}
		\hline
		控温算法 & 优点 & 缺点 \\
		\hline
		PID 控制 & 简单、易实现 & 性能有限，难以处理复杂系统 \\
		\hline
		模糊控制 & 不需要精确模型，适应性强 & 规则设计复杂，性能依赖于规则质量 \\
		\hline
		模型预测控制 & 最优性能，多变量处理 & 计算量大，依赖系统模型 \\
		\hline
		自适应控制 & 参数自调整，适应性强 & 实现复杂，收敛性问题 \\
		\hline
		神经网络控制 & 非线性处理能力强 & 训练复杂，数据需求大 \\
		\hline
		最优控制 & 理论性能优 & 实现复杂，计算量大 \\
		\hline
	\end{tabular}			
\end{table}

来看看\cref{tab:tab1}。

\clearpage
\textbf{这是图片：}

\begin{figure}[htbp]
	\centering
	\includegraphics[width=0.7\textwidth]{Stimulated Raman transition using c.w. laser beams.png}
	\caption{Stimulated Raman transition using c.w. laser beams}
	\label{fig:Stimulated Raman transition using c.w. laser beams}
\end{figure}

来看看\cref{fig:Stimulated Raman transition using c.w. laser beams}。

\clearpage
\subsubsection{（这是第二小节）星和猎手}

\textbf{这个公式用于展示带编号公式的编号规则：}

\begin{equation}
	V_{\text{oc}} = \alpha \Delta T
\end{equation}

\textbf{接下来将用于展示各种各样的自定义环境：}

\textbf{首先是几种颜色盒子：}

\textbf{这是要点：}

\begin{keypoint}[label=key:firefly]
	这是一段关于流萤的文本。
\end{keypoint}

引用示例：\cref{key:firefly}。

\textbf{这是通用框架，通常不引用：}

\begin{ubox}{约会的小曲}
	若我不曾见过太阳、
	使一颗心免于哀伤。
\end{ubox}

\textbf{这是高亮：}

\highlight{流萤！}

\textbf{这是代码环境：}

\begin{macbox}{</>Python}
	\begin{lstlisting}[style=pythonstyle]
		# 示例代码
		import matplotlib.pyplot as plt
		import numpy as np
		
		# Data for plotting
		t = np.arange(0.0, 2.0, 0.01)
		s = 1 + np.sin(2 * np.pi * t)
		
		fig, ax = plt.subplots()
		ax.plot(t, s)
		
		ax.set(xlabel='time (s)', ylabel='voltage (mV)',
		title='About as simple as it gets, folks')
		ax.grid()
		
		fig.savefig("test.png")
		plt.show()
	\end{lstlisting}
\end{macbox}
